\subsection{Тележки и не только}
%%%%%%%%%%%%%%%%%%%%%%%%%%%%%%%%%%%%%%%%%%%%%%%%%%%%
%%%%%%%%%%%%%%%%%%%%%%%%%%%%%%%%%%%%%%%%%%%%%%%%%%%%
 \z[{\begin{wrapfigure}[5]{r}{0.2\textwidth}%
		\vspace{-12pt}%
		\begin{flushright}%
 			\includegraphics[width=0.2\textwidth]{6.4.1}%
		\end{flushright}%
	\end{wrapfigure}%
	}]%
	{На первоначально неподвижной тележке установлены два вертикальных цилиндрических сосуда, соединённых тонкой трубкой. Площадь сечения каждого сосуда $S$, расстояние между их осями $l$. Один из сосудов заполнен жидкостью плотности $\rho$. Кран на соединительной трубке открывают. Найдите скорость тележки $u$ в момент времени, когда скорость уровней жидкости равна $v$. Полная масса всей системы $m$.}
 %
{$u = \frac{\rho S l v}{m}$}
%%%%%%%%%%%%%%%%%%%%%%%%%%%%%%%%%%%%%%%%%%%%%%%%%%%%
 \z[{\begin{wrapfigure}[3]{r}{0.2\textwidth}%
		\vspace{-12pt}%
		\begin{flushright}%
 			\includegraphics[width=0.2\textwidth]{6.4.2}%
		\end{flushright}%
	\end{wrapfigure}%
	}]%
	{На гладком полу стоит сосуд, заполненный водой плотности $\rho_0$; объём воды $V_0$. Оказавшийся на дне сосуда жук объёма $V$ и плотности $\rho$ через некоторое время начинает ползти по дну сосуда со скоростью $u$ относительно него. С какой скоростью $v$ станет двигаться сосуд по полу? Массой сосуда пренебречь, уровень воды всё время остаётся горизонтальным.}
 %
{$v = \frac{(\rho - \rho_0)V u}{\rho_0 V_0 + \rho V}$}
%%%%%%%%%%%%%%%%%%%%%%%%%%%%%%%%%%%%%%%%%%%%%%%%%%%% 
 \z{В сосуде, наполненном водой плотности $\rho$, с ускорением $a$ всплывает пузырёк воздуха, объём которого $V$. Найдите силу давления $F$ со стороны сосуда на опору. Масса сосуда вместе с водой равна $m$.}
%
{$F = m g - \rho V a$}
%%%%%%%%%%%%%%%%%%%%%%%%%%%%%%%%%%%%%%%%%%%%%%%%%%%% 

 \z[{\begin{wrapfigure}[11]{r}{0.15\textwidth}%
		\vspace{-10pt}%
		\begin{flushright}%
 			\includegraphics[width=0.15\textwidth]{6.4.4}%
		\end{flushright}%
	\end{wrapfigure}%
	}]%
	{Обезьяна массы $m$ уравновешена противовесом на блоке $A$. Блок $A$ уравновешен грузом массы $2m$ на блоке $B$. Система неподвижна. Как будет двигаться груз, если обезьяна начнёт равномерно выбирать верёвку со скоростью $u$ относительно себя? Массой блоков и трением пренебречь.}
 %
{$u/4$  вверх}
%%%%%%%%%%%%%%%%%%%%%%%%%%%%%%%%%%%%%%%%%%%%%%%%%%%%
 \z{Платформа массы $M$ в момент времени $t = 0$ начинает двигаться по горизонтальным гладким рельсам под действием постоянной силы тяги $F$, направленной горизонтально. Из неподвижного бункера в момент времени $t = 0$ сверху на неё начинает высыпаться песок. Скорость погрузки постоянна и равна $\mu$ (кг/с). Получите зависимость от времени скорости $v(t)$ и ускорения $a(t)$ платформы в процессе погрузки. Начальная скорость платформы равна нулю.}
%
{$v(t) = \frac{F t}{m_0 \left(1 + \frac{\mu t}{m_0}\right)};$~$a = \frac{F}{m_0 \left(1 + \frac{\mu t}{m_0}\right)^2}$}
